% !TeX root = main_fps.tex
\section{Competing Models}


\subsection{Discrete Element Methods}

Discrete Element Methods (DEM) represent particles as
individual interacting entities whose motion is governed
by contact interactions with neighboring particles and
boundaries. Unlike molecular dynamics methodologies,
which commonly derive forces from potential-energy
functions, DEM formulations typically generate forces
through prescribed contact laws that relate particle
deformation, overlap, relative motion, and material
properties to interaction forces.

A defining characteristic of DEM is the use of contact
models to represent particle interactions. Normal and
tangential forces are frequently computed using
relationships that incorporate stiffness, damping,
friction, cohesion, rolling resistance, or other
material-dependent parameters. These models allow DEM
frameworks to reproduce a wide range of particle
behaviors and have contributed to their widespread use
in the study of granular materials, powders, bulk solids,
and particulate processes.

Although numerous DEM formulations have been proposed,
they generally share a common interaction philosophy.
Particle contacts are identified, a prescribed contact
law is evaluated, forces are generated from the contact
model, and the resulting forces are integrated to produce
particle motion. The resulting causal chain may be
summarized as:

\begin{verbatim}
Contact Law
↓
Force
↓
Acceleration
↓
Motion
\end{verbatim}

Consequently, the behavior of a DEM simulation is closely
linked to the assumptions embedded within the selected
contact model. The choice of stiffness relationships,
friction laws, damping models, and other constitutive
parameters plays a central role in determining the
resulting particle dynamics.

```
```
